\part{Paarungen und Flüsse}

\chapter{Paarungen}

\section{Die Zuteilungsfrage}

\begin{merksatz}
\boxtitel{Vier Ärztinnen, vier Rotationsstellen}
An einer Klinik sind vier Rotationsstellen zu besetzen. Jede Ärztin hat angegeben, für welche Stellen sie die Voraussetzungen mitbringt.

\begin{center}
{\small
\renewcommand{\arraystretch}{1.45}
\begin{tabular}{@{}P{2.0cm}P{5.2cm}@{}}
\toprule
\textbf{Ärztin} & \textbf{Mögliche Stellen}\\
\midrule
$A$ & Intensiv, Notaufnahme\\
$B$ & Intensiv\\
$C$ & Intensiv, Ambulanz\\
$D$ & Notaufnahme, Ambulanz, Sonografie\\
\bottomrule
\end{tabular}}
\end{center}

\textbf{Die Frage:} Lassen sich alle vier gleichzeitig auf verschiedene Stellen setzen?
\end{merksatz}

Man kann probieren. Aber die eigentlich interessante Frage lautet: Woran erkennt man, dass es \kursiv{nicht} geht --- ohne alle Möglichkeiten durchzugehen?

Die Antwort gibt der Satz von Hall, und sie ist von einer Klarheit, die man nicht erwartet.

\begin{definition}
\boxtitelm{Definition 11.1 --- Paarung}
Sei $G = (V,E)$ ein Graph. Eine \Begriff{Paarung} (englisch \kursiv{matching}) ist eine Kantenmenge $M \subseteq E$, in der keine zwei Kanten eine gemeinsame Ecke haben.

\begin{itemize}[leftmargin=1.4em, itemsep=3pt]
\item Eine Ecke heißt \Begriff{gesättigt} durch $M$, wenn eine Kante aus $M$ sie berührt; sonst \Begriff{frei}.
\item $M$ heißt \Begriff{perfekt}, wenn jede Ecke gesättigt ist.
\item $M$ heißt \Begriff{größte Paarung}, wenn es keine Paarung mit mehr Kanten gibt.
\end{itemize}
\end{definition}

\begin{warnung}
\boxtitelr{\enquote{Maximal} und \enquote{größt} sind nicht dasselbe}
Eine Paarung heißt \Begriff{maximal}, wenn man keine Kante mehr hinzufügen kann. Sie heißt \Begriff{größt}, wenn keine andere Paarung mehr Kanten hat. Das erste folgt aus dem zweiten, aber nicht umgekehrt.

\begin{center}
\begin{tikzpicture}[scale=1.25]
\begin{scope}
\node[knoten] (a) at (0,0) {};
\node[knoten] (b) at (1,0) {};
\node[knoten] (c) at (2,0) {};
\node[knoten] (d) at (3,0) {};
\draw[kanted, line width=2pt] (b)--(c);
\draw[kante] (a)--(b); \draw[kante] (c)--(d);
\node[font=\scriptsize\itshape, color=akzent] at (1.5,-0.55) {maximal: eine Kante};
\end{scope}
\begin{scope}[yshift=-1.5cm]
\node[knoten] (e) at (0,0) {};
\node[knoten] (f) at (1,0) {};
\node[knoten] (g) at (2,0) {};
\node[knoten] (h) at (3,0) {};
\draw[kanted, line width=2pt] (e)--(f);
\draw[kanted, line width=2pt] (g)--(h);
\draw[kante] (f)--(g);
\node[font=\scriptsize\itshape, color=gruen] at (1.5,-0.55) {größt: zwei Kanten};
\end{scope}
\end{tikzpicture}
\end{center}

Im oberen Bild lässt sich keine Kante mehr hinzufügen --- alle übrigen berühren $b$ oder $c$. Trotzdem gibt es eine bessere Paarung.

\begin{center}
\begin{minipage}{0.88\textwidth}
\kursiv{Diese Unterscheidung erklärt, warum ein gieriges Verfahren hier nicht genügt: Es liefert eine maximale, aber nicht notwendig eine größte Paarung. Der nächste Satz sagt, woran man den Unterschied erkennt.}
\end{minipage}
\end{center}
\end{warnung}

\section{Vergrößernde Wege}

\begin{definition}
\boxtitelm{Definition 11.2 --- alternierender und vergrößernder Weg}
Sei $M$ eine Paarung.
\begin{itemize}[leftmargin=1.4em, itemsep=3pt]
\item Ein Weg heißt \Begriff{alternierend}, wenn seine Kanten abwechselnd zu $M$ und nicht zu $M$ gehören.
\item Ein alternierender Weg heißt \Begriff{vergrößernd}, wenn seine \kursiv{beiden} Endecken frei sind.
\end{itemize}
\end{definition}

\begin{merksatz}
\boxtitel{Warum der Name}
\begin{center}
\begin{tikzpicture}[scale=1.2]
\node[knoten] (v0) at (0,0) {};
\node[knoten] (v1) at (1.1,0) {};
\node[knoten] (v2) at (2.2,0) {};
\node[knoten] (v3) at (3.3,0) {};
\node[knoten] (v4) at (4.4,0) {};
\draw[kante] (v0)--(v1);
\draw[kanted, line width=2pt] (v1)--(v2);
\draw[kante] (v2)--(v3);
\draw[kanted, line width=2pt] (v3)--(v4);
\node[font=\scriptsize, color=grau] at (0,-0.4) {frei};
\node[font=\scriptsize, color=grau] at (4.4,-0.4) {};
\end{tikzpicture}
\end{center}

Ein vergrößernder Weg beginnt und endet mit einer Nicht-Paarungskante. Er hat also eine ungerade Länge, und die Nicht-Paarungskanten sind um eine zahlreicher als die Paarungskanten.

\schritt{Der Tausch}
Vertauscht man auf dem Weg die Rollen --- jede Paarungskante wird gelöst, jede andere aufgenommen ---, so entsteht wieder eine Paarung, und zwar mit \kursiv{einer Kante mehr}.

\schritt{Warum das wieder eine Paarung ist}
Die inneren Ecken des Weges waren vorher durch je eine Paarungskante gesättigt und sind es nachher wieder --- nur durch eine andere. Die beiden Endecken waren frei und sind nun gesättigt. Ecken außerhalb des Weges sind unberührt.
\end{merksatz}

\begin{merksatz}
\boxtitel{Satz 11.3 (Berge 1957)}
Eine Paarung $M$ ist genau dann eine größte Paarung, wenn es keinen bezüglich $M$ vergrößernden Weg gibt.
\end{merksatz}

\begin{merksatz}
\boxtitel{Beweis}
\schritt{Richtung 1 --- gibt es einen, so ist $M$ nicht größt}
Das ist der Tausch oben: Er liefert eine Paarung mit mehr Kanten.

\schritt{Richtung 2 --- ist $M$ nicht größt, so gibt es einen}
Sei $M'$ eine Paarung mit $|M'| > |M|$.

\schritt{Die symmetrische Differenz}
Wir betrachten
\[ D = (M \setminus M') \cup (M' \setminus M) , \]
also die Kanten, die in genau einer der beiden Paarungen liegen. Sei $H$ der Graph mit Eckenmenge $V$ und Kantenmenge $D$.

\schritt{Welche Grade in $H$ vorkommen}
Eine Ecke $v$ liegt in $M$ auf höchstens einer Kante und in $M'$ auf höchstens einer. Also berühren sie in $H$ höchstens zwei Kanten:
\[ \deg_H(v) \le 2 . \]

\schritt{Wie $H$ aussieht}
Ein Graph mit lauter Graden $\le 2$ zerfällt in Wege und Kreise. (Begründung: Jede Komponente ist zusammenhängend mit Maximalgrad $2$; hat sie eine Ecke vom Grad $\le 1$, so ist sie ein Weg, sonst sind alle Grade genau $2$ und sie ist ein Kreis --- das folgt aus Satz 4.2 und Satz 2.8.)

\schritt{Was auf diesen Wegen und Kreisen steht}
In $H$ liegen keine zwei benachbarten Kanten aus derselben Paarung --- sonst hätten diese eine gemeinsame Ecke innerhalb von $M$ oder innerhalb von $M'$. Also alternieren die Kanten entlang jedes Weges und jedes Kreises zwischen $M$ und $M'$.

\schritt{Der Zählschluss}
Auf einem Kreis alternierender Kanten liegen gleich viele aus $M$ wie aus $M'$ --- die Länge ist gerade.

Auf einem Weg alternierender Kanten sind die Anzahlen gleich (gerade Länge) oder unterscheiden sich um genau eins (ungerade Länge).

Da insgesamt $|M' \setminus M| > |M \setminus M'|$ gilt --- das folgt aus $|M'| > |M|$, weil die gemeinsamen Kanten auf beiden Seiten gleich zählen ---, muss es eine Komponente geben, in der mehr $M'$-Kanten als $M$-Kanten liegen. Das kann nur ein Weg ungerader Länge sein, der mit $M'$-Kanten beginnt und endet.

\schritt{Dieser Weg ist vergrößernd}
Seine beiden Endecken werden in $H$ nur von einer $M'$-Kante berührt. Läge an einer solchen Endecke eine $M$-Kante, so wäre sie entweder in $D$ (dann hätte die Ecke Grad $2$ in $H$ und der Weg wäre dort nicht zu Ende) oder in $M \cap M'$ --- aber dann wäre die Endecke schon durch diese Kante in $M'$ gesättigt und könnte nicht zusätzlich auf einer $M'$-Kante aus $D$ liegen.

Also sind beide Endecken bezüglich $M$ frei, und der Weg alterniert. Er ist vergrößernd. $\blacksquare$
\end{merksatz}

\section{Der Satz von Hall}

Nun der Hauptsatz für den bipartiten Fall --- die Antwort auf die Zuteilungsfrage.

\begin{merksatz}
\boxtitel{Satz 11.4 (Hall 1935)}
Sei $G$ bipartit mit Bipartition $(A,B)$. Dann gibt es genau dann eine Paarung, die \kursiv{alle} Ecken von $A$ sättigt, wenn
\[ |N(S)| \;\ge\; |S| \qquad \text{für jede Teilmenge } S \subseteq A \]
gilt. Dabei ist $N(S) = \bigcup_{v \in S} N(v)$ die Menge aller Nachbarn von Ecken aus $S$.
\end{merksatz}

\begin{merksatz}
\boxtitel{Die Bedingung in Worten}
\begin{center}
\begin{minipage}{0.9\textwidth}
\textbf{Je $k$ Ärztinnen müssen zusammen mindestens $k$ verschiedene Stellen zur Auswahl haben.}
\end{minipage}
\end{center}

Diese Forderung ist offensichtlich \kursiv{notwendig}: Sollen $k$ Ärztinnen auf verschiedene Stellen, so braucht man mindestens $k$ Stellen, und diese müssen für die $k$ Ärztinnen infrage kommen.

Die Leistung des Satzes liegt darin, dass sie auch \kursiv{genügt}. Das ist keineswegs selbstverständlich: Die Bedingung schaut nur auf Teilmengen und Nachbarschaftsgrößen, nicht darauf, wie die Kanten verlaufen.
\end{merksatz}

\begin{merksatz}
\boxtitel{Beweis der notwendigen Richtung}
Sei $M$ eine Paarung, die $A$ sättigt, und sei $S \subseteq A$.

Jede Ecke $a \in S$ ist durch eine Kante $\{a, m(a)\} \in M$ gesättigt. Die Ecke $m(a)$ ist ein Nachbar von $a$, liegt also in $N(S)$.

Verschiedene Ecken $a \neq a'$ liefern verschiedene Partner: Wäre $m(a) = m(a')$, so hätten zwei Kanten von $M$ eine gemeinsame Ecke --- das verbietet Definition 11.1.

Die Zuordnung $a \mapsto m(a)$ ist somit eine injektive Abbildung von $S$ nach $N(S)$, also $|S| \le |N(S)|$. $\blacksquare$
\end{merksatz}

\begin{merksatz}
\boxtitel{Beweis der hinreichenden Richtung --- Induktion über $|A|$}
\schritt{Induktionsanfang $|A| = 1$}
Sei $A = \{a\}$. Die Bedingung für $S = \{a\}$ lautet $|N(a)| \ge 1$, also hat $a$ einen Nachbarn $b$. Die Paarung $\{\{a,b\}\}$ sättigt $A$.

\schritt{Induktionsschritt --- die Fallunterscheidung}
Sei $|A| \ge 2$ und die Behauptung für kleinere $A$ bewiesen. Wir unterscheiden danach, ob die Bedingung irgendwo \kursiv{knapp} erfüllt ist.

\schritt{Fall 1 --- überall Luft}
Angenommen, für jede nichtleere echte Teilmenge $S \subsetneq A$ gilt sogar
\[ |N(S)| \;\ge\; |S| + 1 . \]

Wir wählen eine beliebige Ecke $a \in A$ und einen ihrer Nachbarn $b$ (er existiert nach der Bedingung für $S = \{a\}$). Diese Kante nehmen wir in die Paarung auf.

Nun betrachten wir $G' = G - a - b$ mit $A' = A \setminus \{a\}$.

\kursiv{Erfüllt $G'$ die Hall-Bedingung?} Sei $S \subseteq A'$ nichtleer. In $G'$ fehlt gegenüber $G$ höchstens die Ecke $b$ aus der Nachbarschaft:
\[ |N_{G'}(S)| \;\ge\; |N_G(S)| - 1 \;\ge\; (|S| + 1) - 1 \;=\; |S| . \]
Dabei wurde benutzt, dass $S$ eine nichtleere echte Teilmenge von $A$ ist --- echt, weil $a \notin S$.

Nach Induktionsannahme gibt es eine Paarung $M'$ in $G'$, die $A'$ sättigt. Zusammen mit $\{a,b\}$ sättigt sie ganz $A$.

\schritt{Fall 2 --- irgendwo ist es knapp}
Angenommen, es gibt eine nichtleere echte Teilmenge $S_0 \subsetneq A$ mit
\[ |N(S_0)| = |S_0| . \]

\kursiv{Der erste Teilgraph.} Sei $G_1$ der von $S_0 \cup N(S_0)$ induzierte Teilgraph. Für $S \subseteq S_0$ ist $N_{G_1}(S) = N_G(S)$ --- alle Nachbarn von Ecken aus $S_0$ liegen ja in $N(S_0)$. Die Hall-Bedingung gilt also in $G_1$.

Da $|S_0| < |A|$, liefert die Induktionsannahme eine Paarung $M_1$, die $S_0$ sättigt. Wegen $|N(S_0)| = |S_0|$ sättigt sie sogar ganz $N(S_0)$.

\kursiv{Der zweite Teilgraph.} Sei $G_2$ der von $(A \setminus S_0) \cup (B \setminus N(S_0))$ induzierte Teilgraph, mit $A_2 = A \setminus S_0$.

\kursiv{Erfüllt $G_2$ die Hall-Bedingung?} Sei $T \subseteq A_2$. Wir betrachten $T \cup S_0 \subseteq A$ und wenden die Bedingung in $G$ an:
\[ |N_G(T \cup S_0)| \;\ge\; |T \cup S_0| = |T| + |S_0| , \]
wobei die letzte Gleichheit gilt, weil $T$ und $S_0$ disjunkt sind.

Weiter ist $N_G(T \cup S_0) = N_G(T) \cup N_G(S_0)$, und die Nachbarn von $T$ in $G_2$ sind genau diejenigen Nachbarn in $G$, die nicht in $N(S_0)$ liegen. Also
\[ |N_{G_2}(T)| = |N_G(T \cup S_0)| - |N_G(S_0)| \;\ge\; \bigl(|T| + |S_0|\bigr) - |S_0| = |T| . \]

Da $S_0$ nichtleer ist, gilt $|A_2| < |A|$. Nach Induktionsannahme gibt es eine Paarung $M_2$, die $A_2$ sättigt.

\schritt{Zusammensetzen}
$M_1$ benutzt nur Ecken aus $S_0 \cup N(S_0)$, $M_2$ nur Ecken außerhalb davon. Die Vereinigung $M_1 \cup M_2$ ist daher eine Paarung, und sie sättigt $S_0 \cup A_2 = A$. $\blacksquare$
\end{merksatz}

\begin{merksatz}
\boxtitel{Die Zuteilungsfrage beantwortet}
\schritt{Der Graph}
$A = \{A, B, C, D\}$ (Ärztinnen), $B = \{$Int, Not, Amb, Son$\}$.

\begin{center}
\begin{tikzpicture}[scale=1.25]
\node[knotenb] (a) at (-2.4,0.9) {$A$};
\node[knotenb] (b) at (-0.8,0.9) {$B$};
\node[knotenb] (c) at (0.8,0.9) {$C$};
\node[knotenb] (d) at (2.4,0.9) {$D$};
\node[knotenr] (i) at (-2.4,-0.9) {I};
\node[knotenr] (n) at (-0.8,-0.9) {N};
\node[knotenr] (m) at (0.8,-0.9) {A};
\node[knotenr] (s) at (2.4,-0.9) {S};
\draw[kante] (a)--(i); \draw[kante] (a)--(n);
\draw[kante] (b)--(i);
\draw[kante] (c)--(i); \draw[kante] (c)--(m);
\draw[kante] (d)--(n); \draw[kante] (d)--(m); \draw[kante] (d)--(s);
\end{tikzpicture}
\end{center}

\schritt{Die Hall-Bedingung an der kritischen Stelle prüfen}
Man prüft nicht alle $2^4 = 16$ Teilmengen, sondern sucht gezielt nach knappen Stellen --- also nach Gruppen mit wenigen gemeinsamen Möglichkeiten.

\begin{center}
{\small
\renewcommand{\arraystretch}{1.42}
\begin{tabular}{@{}P{2.0cm}P{2.4cm}P{1.2cm}P{1.4cm}@{}}
\toprule
$S$ & $N(S)$ & $|S|$ & $|N(S)|$\\
\midrule
$\{B\}$ & Int & $1$ & $1$\\
$\{A,B\}$ & Int, Not & $2$ & $2$\\
$\{B,C\}$ & Int, Amb & $2$ & $2$\\
$\{A,B,C\}$ & Int, Not, Amb & $3$ & $3$\\
$\{A,B,C,D\}$ & alle vier & $4$ & $4$\\
\bottomrule
\end{tabular}}
\end{center}

Alle geprüften Teilmengen erfüllen $|N(S)| \ge |S|$; die übrigen erst recht, da sie größere Nachbarschaften haben. \textbf{Die Zuteilung ist möglich.}

\schritt{Eine Zuteilung angeben}
$B \to$ Intensiv (die einzige Möglichkeit), $A \to$ Notaufnahme, $C \to$ Ambulanz, $D \to$ Sonografie.

\schritt{Was passiert, wenn $D$ ausfällt?}
Dann bleiben $A$, $B$, $C$ mit den Möglichkeiten Int, Not / Int / Int, Amb. Prüfung: $|N(\{A,B,C\})| = 3 = |S|$. Es geht noch --- gerade eben.

\schritt{Und wenn zusätzlich die Ambulanz wegfällt?}
Dann hat $C$ nur noch Intensiv, ebenso $B$. Für $S = \{B,C\}$ ist $|N(S)| = 1 < 2$. \textbf{Unmöglich} --- und man kann genau angeben, warum: Zwei Ärztinnen bewerben sich um eine einzige Stelle.

\begin{center}
\begin{minipage}{0.88\textwidth}
\kursiv{Das ist der praktische Wert des Satzes. Er sagt nicht nur \enquote{geht nicht}, sondern liefert die Begründung in Gestalt einer konkreten Gruppe --- ein Argument, das man einem Vorgesetzten vorlegen kann.}
\end{minipage}
\end{center}
\end{merksatz}

\begin{merksatz}
\boxtitel{Folgerung 11.5 --- reguläre bipartite Graphen}
Ist $G$ bipartit und $r$-regulär mit $r \ge 1$ --- also hat jede Ecke genau den Grad $r$ ---, so besitzt $G$ eine perfekte Paarung.

\schritt{Beweis, Schritt 1: die Seiten sind gleich groß}
Wir zählen die Kanten auf zwei Weisen. Jede Kante hat genau ein Ende in $A$ und eines in $B$, also
\[ m = \sum_{a \in A}\deg(a) = r|A| \qquad \text{und} \qquad m = \sum_{b \in B}\deg(b) = r|B| . \]
Wegen $r \ge 1$ folgt $|A| = |B|$.

\schritt{Schritt 2: die Hall-Bedingung}
Sei $S \subseteq A$. Wir zählen die Kanten, die in $S$ beginnen, auf zwei Weisen.

\kursiv{Von $S$ aus:} Es sind genau $r|S|$ Kanten.

\kursiv{Von $N(S)$ aus:} Jede dieser Kanten endet in $N(S)$. Jede Ecke von $N(S)$ hat insgesamt $r$ Kanten, von denen höchstens alle nach $S$ führen können. Also enden höchstens $r|N(S)|$ Kanten dort.

Somit
\[ r|S| \;\le\; r|N(S)| \quad\Longrightarrow\quad |S| \le |N(S)| . \]

\schritt{Schritt 3: Schluss}
Nach Hall gibt es eine Paarung, die $A$ sättigt. Sie hat $|A|$ Kanten und sättigt damit $|A| = |B|$ Ecken in $B$ --- also alle. Sie ist perfekt. $\blacksquare$

\begin{center}
\begin{minipage}{0.88\textwidth}
\kursiv{Diese Folgerung hat eine hübsche Anwendung: Ein Kartenspiel mit $52$ Karten wird in $13$ Stapel zu je $4$ Karten geteilt. Dann kann man stets aus jedem Stapel eine Karte so wählen, dass alle dreizehn Werte --- Zwei bis Ass --- genau einmal vorkommen. Man setze Stapel gegen Werte und verbinde, wenn der Wert im Stapel vorkommt; der Graph ist $4$-regulär.}
\end{minipage}
\end{center}
\end{merksatz}

\section{Der Satz von König}

\begin{definition}
\boxtitelm{Definition 11.6 --- Eckenüberdeckung}
Eine \Begriff{Eckenüberdeckung} ist eine Eckenmenge $C \subseteq V$, so dass jede Kante mindestens eine ihrer Endecken in $C$ hat.
\end{definition}

\begin{merksatz}
\boxtitel{Der einfache Zusammenhang}
Für jeden Graphen gilt: Ist $M$ eine Paarung und $C$ eine Eckenüberdeckung, so ist
\[ |M| \;\le\; |C| . \]

\schritt{Warum}
Jede Kante von $M$ muss von $C$ überdeckt werden, braucht also mindestens eine Ecke aus $C$. Verschiedene Kanten von $M$ haben keine gemeinsame Ecke, brauchen also verschiedene Ecken aus $C$. $\blacksquare$

\begin{center}
\begin{minipage}{0.88\textwidth}
\kursiv{Die Ungleichung gilt immer. Die Frage ist, wann Gleichheit eintritt. Im Dreieck $K_3$ ist die größte Paarung $1$ und die kleinste Überdeckung $2$ --- Gleichheit gilt also nicht allgemein.}
\end{minipage}
\end{center}
\end{merksatz}

\begin{merksatz}
\boxtitel{Satz 11.7 (Kőnig 1931, Egerváry 1931)}
In einem bipartiten Graphen ist die Größe einer größten Paarung gleich der Größe einer kleinsten Eckenüberdeckung.
\end{merksatz}

\begin{merksatz}
\boxtitel{Beweis --- die Konstruktion}
Sei $M$ eine größte Paarung, Bipartition $(A,B)$.

\schritt{Schritt 1 --- Die erreichbaren Ecken}
Sei $U \subseteq A$ die Menge der bezüglich $M$ freien Ecken in $A$.

Sei $Z$ die Menge aller Ecken, die von einer Ecke aus $U$ über einen \kursiv{alternierenden} Weg erreichbar sind (einschließlich $U$ selbst, über Wege der Länge $0$). Dabei beginnt jeder solche Weg mit einer Nicht-Paarungskante, weil die Startecken frei sind.

Wir setzen
\[ Z_A = Z \cap A, \qquad Z_B = Z \cap B . \]

\schritt{Schritt 2 --- Die vorgeschlagene Überdeckung}
\[ C = (A \setminus Z_A) \;\cup\; Z_B . \]

\begin{center}
\begin{tikzpicture}[scale=1.15]
\draw[kante, grau!45] (-1.6,0.9) ellipse (1.5 and 0.55);
\draw[kante, grau!45] (1.6,0.9) ellipse (1.5 and 0.55);
\draw[kante, grau!45] (-1.6,-0.9) ellipse (1.5 and 0.55);
\draw[kante, grau!45] (1.6,-0.9) ellipse (1.5 and 0.55);
\node[font=\scriptsize, color=grau] at (-1.6,0.9) {$Z_A$};
\node[font=\scriptsize, color=akzent] at (1.6,0.9) {$A \setminus Z_A$};
\node[font=\scriptsize, color=akzent] at (-1.6,-0.9) {$Z_B$};
\node[font=\scriptsize, color=grau] at (1.6,-0.9) {$B \setminus Z_B$};
\node[font=\scriptsize\itshape, color=akzent] at (0,-1.75) {rot: die Überdeckung $C$};
\end{tikzpicture}
\end{center}

\schritt{Schritt 3 --- $C$ ist wirklich eine Überdeckung}
Angenommen, eine Kante $\{a,b\}$ mit $a \in A$, $b \in B$ wäre nicht überdeckt. Dann ist $a \notin A \setminus Z_A$, also $a \in Z_A$; und $b \notin Z_B$.

$a \in Z_A$ heißt: Es gibt einen alternierenden Weg $P$ von einer freien Ecke aus $U$ nach $a$. Dieser Weg endet in $a \in A$; da er mit einer Nicht-Paarungskante beginnt und alterniert, endet er mit einer \kursiv{Paarungs}kante (oder hat Länge $0$, falls $a$ selbst frei ist).

In beiden Fällen können wir $P$ um die Kante $\{a,b\}$ verlängern --- sie ist keine Paarungskante, denn wäre sie eine, so wäre $b$ über $P$ erreichbar und läge in $Z_B$. Der verlängerte Weg ist alternierend und erreicht $b$. Also $b \in Z_B$ --- Widerspruch. \checkmark

\schritt{Schritt 4 --- $|C| \le |M|$}
Wir zeigen: Jede Ecke von $C$ ist durch $M$ gesättigt, und verschiedene Ecken von $C$ werden durch \kursiv{verschiedene} Kanten von $M$ gesättigt.

\kursiv{Ecken in $A \setminus Z_A$ sind gesättigt.} Wäre eine solche Ecke frei, so läge sie in $U \subseteq Z_A$ --- Widerspruch.

\kursiv{Ecken in $Z_B$ sind gesättigt.} Wäre $b \in Z_B$ frei, so wäre der alternierende Weg von $U$ nach $b$ ein \kursiv{vergrößernder} Weg: Er beginnt an einer freien Ecke in $A$ und endet an einer freien Ecke in $B$. Nach Satz 11.3 wäre $M$ dann nicht größt --- Widerspruch.

\kursiv{Keine Paarungskante hat beide Enden in $C$.} Sei $\{a,b\} \in M$ mit $b \in Z_B$. Der alternierende Weg von $U$ nach $b$ endet mit einer Nicht-Paarungskante (er beginnt mit einer solchen und alterniert, und $b \in B$). Verlängern wir ihn um die Paarungskante $\{b,a\}$, so erreichen wir $a$ auf alternierendem Weg. Also $a \in Z_A$, und damit $a \notin A \setminus Z_A$.

Somit liefert jede Ecke von $C$ eine eigene Kante von $M$, und
\[ |C| \le |M| . \]

\schritt{Schritt 5 --- Schluss}
Zusammen mit der allgemeinen Ungleichung $|M| \le |C'|$ für jede Überdeckung $C'$ folgt: $C$ ist eine kleinste Überdeckung und $|C| = |M|$. $\blacksquare$
\end{merksatz}

\begin{merksatz}
\boxtitel{Satz 11.8 (Tutte 1947) --- der allgemeine Fall}
Ein Graph $G$ besitzt genau dann eine perfekte Paarung, wenn für jede Eckenmenge $S \subseteq V$ gilt:
\[ \text{Die Anzahl der Komponenten \kursiv{ungerader} Eckenzahl von } G - S \ \text{ ist } \le |S| . \]

\begin{center}
\begin{minipage}{0.88\textwidth}
\kursiv{Auch diesen Satz beweise ich nicht --- wie Kuratowski und Brooks. Die notwendige Richtung ist leicht (jede ungerade Komponente muss eine Ecke nach $S$ hinausheiraten), die hinreichende verlangt eine sorgfältige Induktion über mehrere Seiten. Diestel\zit{7}, Kapitel~2 führt sie aus.}
\end{minipage}
\end{center}

\schritt{Warum der Satz wichtig ist}
Hall gilt nur bipartit. Tutte gilt allgemein --- und liefert sofort, dass der Petersen-Graph eine perfekte Paarung hat: Man müsste ein $S$ finden, das den Satz verletzt, und für $3$-reguläre, zusammenhängende Graphen ohne Brücke lässt sich zeigen, dass keines existiert (Satz von Petersen 1891).
\end{merksatz}

\section{Übungsaufgaben}

\begin{aufgabe}
\boxtitel{Aufgabe 11.1 --- Hall anwenden}
Sechs Studierende sollen sechs Prüfungsthemen zugeteilt werden. Die Wünsche:
\begin{center}
{\small
\renewcommand{\arraystretch}{1.38}
\begin{tabular}{@{}P{1.4cm}P{2.6cm}P{1.4cm}P{2.6cm}@{}}
\toprule
$S_1$ & $T_1, T_2$ & $S_4$ & $T_3, T_4$\\
$S_2$ & $T_1, T_2$ & $S_5$ & $T_3, T_4$\\
$S_3$ & $T_1, T_2$ & $S_6$ & $T_5, T_6$\\
\bottomrule
\end{tabular}}
\end{center}
Ist eine Zuteilung möglich?

\loesung
\schritt{Schritt 1 --- Nach knappen Stellen suchen}
Die drei Studierenden $S_1$, $S_2$, $S_3$ haben zusammen nur die Themen $T_1$ und $T_2$ zur Auswahl.

\schritt{Schritt 2 --- Die Bedingung prüfen}
Für $S = \{S_1, S_2, S_3\}$:
\[ N(S) = \{T_1, T_2\}, \qquad |N(S)| = 2 \;<\; 3 = |S| . \]

\schritt{Antwort}
\textbf{Nein.} Die Hall-Bedingung ist verletzt, also gibt es keine Zuteilung.

\schritt{Die Begründung in Worten}
Drei Studierende bewerben sich um zwei Themen. Einer von ihnen muss leer ausgehen --- unabhängig davon, wie man die übrigen verteilt.

\schritt{Schritt 3 --- Die kleinste Reparatur}
Gibt man einem der drei ein zusätzliches Thema, etwa $S_3$ auch $T_5$, so prüfen wir neu:
\begin{center}
{\small
\renewcommand{\arraystretch}{1.38}
\begin{tabular}{@{}P{2.6cm}P{2.6cm}P{1.0cm}P{1.2cm}@{}}
\toprule
$S$ & $N(S)$ & $|S|$ & $|N(S)|$\\
\midrule
$\{S_1,S_2\}$ & $T_1,T_2$ & $2$ & $2$\\
$\{S_1,S_2,S_3\}$ & $T_1,T_2,T_5$ & $3$ & $3$\\
$\{S_4,S_5\}$ & $T_3,T_4$ & $2$ & $2$\\
$\{S_3,S_6\}$ & $T_1,T_2,T_5,T_6$ & $2$ & $4$\\
alle sechs & alle sechs & $6$ & $6$\\
\bottomrule
\end{tabular}}
\end{center}
Kein Verstoß mehr. Eine Zuteilung: $S_1 \to T_1$, $S_2 \to T_2$, $S_3 \to T_5$, $S_4 \to T_3$, $S_5 \to T_4$, $S_6 \to T_6$. \checkmark
\end{aufgabe}

\begin{aufgabe}
\boxtitel{Aufgabe 11.2 --- König nachrechnen}
Bestimmen Sie im abgebildeten bipartiten Graphen eine größte Paarung und eine kleinste Eckenüberdeckung.

\begin{center}
\begin{tikzpicture}[scale=1.25]
\node[knotenb] (a1) at (-1.8,0.9) {$a_1$};
\node[knotenb] (a2) at (-0.6,0.9) {$a_2$};
\node[knotenb] (a3) at (0.6,0.9) {$a_3$};
\node[knotenb] (a4) at (1.8,0.9) {$a_4$};
\node[knotenr] (b1) at (-1.8,-0.9) {$b_1$};
\node[knotenr] (b2) at (-0.6,-0.9) {$b_2$};
\node[knotenr] (b3) at (0.6,-0.9) {$b_3$};
\node[knotenr] (b4) at (1.8,-0.9) {$b_4$};
\draw[kante] (a1)--(b1); \draw[kante] (a1)--(b2);
\draw[kante] (a2)--(b1);
\draw[kante] (a3)--(b1); \draw[kante] (a3)--(b3);
\draw[kante] (a4)--(b3); \draw[kante] (a4)--(b4);
\end{tikzpicture}
\end{center}

\loesung
\schritt{Schritt 1 --- Eine Paarung suchen}
Wir versuchen es gierig: $a_1b_1$, dann $a_3b_3$, dann $a_4b_4$. Die Ecke $a_2$ hat nur $b_1$ als Nachbarn, und $b_1$ ist belegt.

Paarung $M = \{a_1b_1,\ a_3b_3,\ a_4b_4\}$ mit $|M| = 3$.

\schritt{Schritt 2 --- Gibt es einen vergrößernden Weg?}
Freie Ecken in $A$: nur $a_2$. Freie Ecken in $B$: nur $b_2$.

Ein vergrößernder Weg müsste von $a_2$ nach $b_2$ laufen und alternieren. Von $a_2$ führt nur die Kante zu $b_1$ (keine Paarungskante \checkmark). Von $b_1$ geht es über die Paarungskante zu $a_1$. Von $a_1$ über eine Nicht-Paarungskante --- das ist $a_1b_2$. Und $b_2$ ist frei.

\textbf{Gefunden:} $a_2 - b_1 - a_1 - b_2$.

\schritt{Schritt 3 --- Tauschen}
Wir nehmen $a_2b_1$ und $a_1b_2$ auf, entfernen $a_1b_1$:
\[ M' = \{a_2b_1,\ a_1b_2,\ a_3b_3,\ a_4b_4\}, \qquad |M'| = 4 . \]

\schritt{Schritt 4 --- Ist $M'$ größt?}
$|A| = 4$ und alle vier Ecken sind gesättigt. Mehr als $4$ Kanten kann keine Paarung haben, denn jede Kante verbraucht eine Ecke aus $A$. Also ja.

\schritt{Schritt 5 --- Die Überdeckung}
Nach König muss die kleinste Überdeckung ebenfalls Größe $4$ haben. Eine solche: $C = \{a_1, a_2, a_3, a_4\}$ --- jede Kante hat ihr $A$-Ende darin.

\schritt{Probe}
Kleiner geht es nicht: Nach dem einfachen Zusammenhang ist $|C| \ge |M'| = 4$ für jede Überdeckung. \checkmark

\begin{center}
\begin{minipage}{0.86\textwidth}
\kursiv{Beachten Sie, dass die gierige Paarung aus Schritt 1 maximal, aber nicht größt war --- genau der Fall aus dem Warnkasten zu Beginn. Der vergrößernde Weg hat sie repariert.}
\end{minipage}
\end{center}
\end{aufgabe}

\begin{aufgabe}
\boxtitel{Aufgabe 11.3 --- Ein Satz über Paarungen in Bäumen}
Zeigen Sie: Ein Baum hat höchstens eine perfekte Paarung.

\loesung
\schritt{Schritt 1 --- Annahme}
Seien $M_1 \neq M_2$ zwei perfekte Paarungen eines Baumes $T$.

\schritt{Schritt 2 --- Die symmetrische Differenz}
Wie im Beweis von Satz 11.3 betrachten wir $D = (M_1 \setminus M_2) \cup (M_2 \setminus M_1)$. Da $M_1 \neq M_2$, ist $D$ nicht leer.

\schritt{Schritt 3 --- Grade in $D$}
Jede Ecke liegt in $M_1$ auf genau einer Kante (perfekt!) und in $M_2$ ebenfalls auf genau einer.

Ist eine Ecke $v$ überhaupt an $D$ beteiligt, so sind ihre beiden Kanten verschieden --- sonst läge die gemeinsame Kante in $M_1 \cap M_2$ und käme in $D$ nicht vor. Beide gehören dann zu $D$, also
\[ \deg_D(v) = 2 \quad \text{für jede von } D \text{ berührte Ecke.} \]

\schritt{Schritt 4 --- Der Widerspruch}
Ein nichtleerer Graph, in dem jede berührte Ecke Grad $2$ hat, enthält einen Kreis. (Das ist der Hilfssatz aus dem Beweis von Satz 6.4.)

Ein Baum ist kreisfrei. Widerspruch.

\schritt{Schritt 5 --- Schluss}
Also $M_1 = M_2$. $\blacksquare$

\schritt{Nachbemerkung --- die Existenz ist eine andere Frage}
Der Satz sagt nur \enquote{höchstens eine}. Viele Bäume haben gar keine: Der Stern $K_{1,3}$ etwa hat vier Ecken, aber jede Paarung kann höchstens eine Kante enthalten --- alle Kanten treffen sich in der Mitte.
\end{aufgabe}

\begin{geschichte}
\boxtitelg{Ein Satz mit vier Vätern und einem irreführenden Namen}
Der Satz von Hall heißt im Englischen \kursiv{marriage theorem} --- Heiratssatz. Die Einkleidung stammt von \Person{Hermann Weyl} (1949): $n$ Damen und $n$ Herren, jede Dame kennt einige Herren, wann lassen sich alle verheiraten?

\begin{center}
{\small
\renewcommand{\arraystretch}{1.42}
\begin{tabular}{@{}P{1.4cm}P{2.4cm}P{4.0cm}@{}}
\toprule
\textbf{1912} & Frobenius & äquivalente Determinantenaussage\\
\textbf{1931} & Kőnig, Egerváry & Paarung $=$ Überdeckung, bipartit\\
\textbf{1935} & Hall & die Bedingung in heutiger Form\\
\textbf{1949} & Weyl & die Einkleidung als Heiratsproblem\\
\bottomrule
\end{tabular}}
\end{center}

\Person{Philip Hall} stellte den Satz nicht als Zuteilungsproblem, sondern als Frage über \kursiv{Repräsentantensysteme}: Gegeben endlich viele Mengen --- wann kann man aus jeder ein Element so wählen, dass alle Gewählten verschieden sind?

\begin{center}
\begin{minipage}{0.9\textwidth}
\kursiv{Die Heiratseinkleidung hat dem Satz zu Bekanntheit verholfen und ihm zugleich geschadet. Sie legt nahe, es gehe um Paare und Symmetrie --- tatsächlich ist die Aussage unsymmetrisch: Sie sättigt \kursiv{eine} Seite und sagt über die andere nichts. In der Zuteilungsfrage dieses Kapitels ist genau das der Punkt.}
\end{minipage}
\end{center}
\end{geschichte}

\begin{forschung}
\boxtitelp{Von Edmonds zur Nierenlebendspende}
\begin{center}
{\small
\renewcommand{\arraystretch}{1.45}
\begin{tabular}{@{}P{1.4cm}P{2.6cm}P{3.8cm}@{}}
\toprule
\textbf{1965} & Edmonds & Verfahren für größte Paarungen in beliebigen Graphen, $\Oh(n^3)$\\
\textbf{1962} & Gale, Shapley & stabile Paarungen mit Präferenzen\\
\textbf{2012} & --- & Wirtschaftsnobelpreis an Shapley und Roth\\
\textbf{seit 2004} & Roth u.\,a. & Überkreuz-Nierenspende\\
\bottomrule
\end{tabular}}
\end{center}

\Person{Jack Edmonds} löste 1965 den allgemeinen Fall. Der Kern seines Verfahrens ist der Umgang mit ungeraden Kreisen --- den \enquote{Blüten} ---, an denen die einfache Suche nach vergrößernden Wegen scheitert. In bipartiten Graphen gibt es solche Kreise nicht, weshalb dort das Verfahren aus Satz 11.3 direkt funktioniert.

Der praktisch folgenreichste Zweig betrifft die Nierenlebendspende. Ein Spender ist oft mit seinem Angehörigen nicht verträglich, aber mit einem fremden Empfänger --- und umgekehrt. Man baut den Graphen der Verträglichkeiten und sucht eine größte Paarung oder, bei längeren Ketten, größte Kreiszerlegungen.

\begin{center}
\begin{minipage}{0.88\textwidth}
\kursiv{Programme dieser Art laufen seit 2004 in den USA und seit 2017 in mehreren europäischen Ländern; die längste dokumentierte Kette umfasste über dreißig Transplantationen. Die zugrundeliegende Mathematik ist die dieses Kapitels --- mit dem Zusatz, dass Kreise fester Länge zugelassen werden müssen, was das Problem NP-schwer macht.}
\end{minipage}
\end{center}
\end{forschung}

\begin{merksatz}
\boxtitel{Was Sie jetzt können}
\begin{center}
\renewcommand{\arraystretch}{1.5}
\begin{tabular}{@{}P{3.0cm}P{5.6cm}@{}}
\toprule
\textbf{Satz 11.3} & Berge: größt $\iff$ kein vergrößernder Weg\\
\textbf{Satz 11.4} & Hall: $|N(S)| \ge |S|$ für alle $S \subseteq A$\\
\textbf{Folgerung 11.5} & Reguläre bipartite Graphen haben perfekte Paarungen\\
\textbf{Satz 11.7} & Kőnig: Paarung $=$ Überdeckung im bipartiten Fall\\
\textbf{Satz 11.8} & Tutte: der allgemeine Fall\\
\bottomrule
\end{tabular}
\end{center}

\begin{center}
\begin{minipage}{0.88\textwidth}
\kursiv{Im nächsten Kapitel zeigt sich, dass Hall und König Sonderfälle eines einzigen, umfassenderen Satzes sind --- und dass derselbe Satz auch die Frage nach Engpässen aus Kapitel~2 beantwortet.}
\end{minipage}
\end{center}
\end{merksatz}
